%!TEX root = ../../../exa-ma-deliverables.tex
\VThreeApplicationDraft{app-cgal-autorefine}
{Mesh autorefinement}
{V. Chabannes; C. Prud'homme; P. Alliez}
{WP1}{mini-app}
{Test the parallel performance of triangle-mesh autorefinement.}
{CGAL, Intel TBB}{benchmark-ready}

\paragraph{Proposed benchmark contract.}
\texttt{BM-CGAL-AUTOREFINE-01} defines a single-node thread-scaling benchmark measuring time-to-solution for the CGAL mesh autorefinement algorithm.
At the inspected commitAt inspected commit \href{https://github.com/numpex/apps-cgal/commit/774f40bb6bf2d42f08468284e55ce7e93d446b4b}{\texttt{774f40b}},
the application repository contains the autorefinement minimal example, input meshes and JSON benchmark and plot configurations under \texttt{src/app-cgal-autorefine/}.

The benchmark has been executed on a single node of the Gaya cluster, for mesh sizes ranging from 100 to 1000 points per grid side and for up to 64 threads.
The resulting measurements are reported in this specification.

A benchmark result is considered valid when the autorefinement application completes successfully and produces the expected autorefined output mesh.
Execution time is measured by the application's internal timing instrumentation for the autorefinement operation. Each configuration is executed once, with the requested number of threads allocated to a single node.


\subsection{Description of the benchmark}

This benchmark evaluates the performance of the mesh autorefinement algorithm provided by CGAL.
The routine takes two grid meshes and merges them, then it refines it so that all triangles and edges are correctly connected.

Input meshes and resulting mesh can be seen on \Cref{fig:autorefine_io}.



\subsection{Benchmarking methodology}
\subsubsection{Benchmarking tools used}

This benchmark was performed on one node of the gaya supercomputer, using feelpp.benchmaring (v4.0.0).
Performance measurements were collected using internal timing instrumentation.

Gaya nodes are equipped with AMD EPYC 7713 64 core processors and 512 GB of RAM.
For these performance measurements, the execution environment was strictly configured to allocate exactly one thread per task. Additionally, each benchmark task was executed in a single pass with no repetitions.

\subsection{Input/Output Dataset Description}

\subsubsection{Input Data:}

\begin{itemize}
\item \textbf{Input Meshes:} Initial two flat grid triangle meshes of varying sizes of side lengths ranging from 100 to 1000 points.
\item \textbf{Setup:} Utilizes the CGAL framework’s mesh refinement functions, with shared-memory multithreading enabled via Intel TBB (Threading Building Blocks).
\item \textbf{Execution Environment:}  Single-node execution to measure shared-memory performance and potential thread contention.
\end{itemize}


\begin{figure}[h!]
  \centering
  \begin{minipage}{0.48\textwidth}
    \centering
    \includegraphics[width=\textwidth]{cgal/autoref_input_meshes.png}
    \subcaption{Input Meshes}
  \end{minipage}\hfill
  \begin{minipage}{0.48\textwidth}
    \centering
    \includegraphics[width=\textwidth]{cgal/autoref_output_mesh.png}
    \subcaption{Output Mesh}
  \end{minipage}
  \caption{Input and Output Meshes for the Autorefinement Benchmark}
  \label{fig:autorefine_io}
\end{figure}



\subsubsection{Output Data:}

The primary monitored quantity is the total execution time of the mesh autorefinement algorithm. These timings are recorded and analyzed to compute the parallel speedup and evaluate the algorithmic time complexity as the problem size scales.

\subsection{Results summary}
\ApplicationBenchmarkResults{app-cgal-autorefine}
